\subsection{Model, controller and optimization}
SCONE's planar \texttt{Human0914} model (OpenSim 3 engine \cite{delp2007}) has 9 degrees of freedom, seven Hill-type muscles per leg (hamstrings, gluteus maximus, iliopsoas, vasti, gastrocnemius, soleus, TA) and compliant foot-ground contact, with a total mass of \SI{75.2}{kg}. It is driven by the Geyer and Herr reflex controller \cite{geyer2010} as implemented in the SCONE tutorials. The objective is the tutorial gait measure with the minimum speed raised to \SI{1.2}{m/s} (speeds more than 5\% below it are penalized), the Wang et al.\ metabolic cost of transport \cite{wang2012} with weight 0.1, and soft joint limits. Each evaluation simulates \SI{10}{s}; a run that ends early because the centre of mass drops is a fall. Optimizations use CMA-ES \cite{hansen2016} and are called converged when the best cost improved by less than 1\% over the last 40 generations. Every condition has three seeds (1, 2, 3); all statistics are mean and standard deviation (SD) over the seeds that walk, and falls are listed separately.

The healthy model was optimized from the tutorial result. Its stride-averaged right-leg curves (strides after the first two) are compared with the 1.2~m/s level treadmill walking of the Camargo et al.\ dataset \cite{camargo2021} (22 able-bodied adults, 568 strides) by RMSE and Pearson $r$ per curve.

\subsection{Drop-foot model and clinical metrics}
Drop foot is a scaled maximum isometric force of \texttt{tib\_ant\_r}: 50, 25 and 10\% of healthy. The \emph{immediate} effect uses the healthy controller unchanged (written once per side); the weakest strength that still walks was found by bisection. The \emph{adapted} gait re-optimizes an asymmetric controller (independent reflex parameters per leg) for 150 generations from the healthy result; seeds that had not converged by then were continued from their best parameters until they did. Metrics are computed per right stride and averaged:
minimum toe clearance (height of the \texttt{toes} origin, the metatarsophalangeal joint, over the middle half of swing; the primary metric);
ankle angle at initial contact;
peak swing dorsiflexion;
a foot-slap index, the peak plantarflexion velocity in the first 15\% of the cycle;
steppage, the change of peak swing hip and knee flexion from healthy;
peak ankle push-off power in stance;
step-length and stance-time symmetry indices \cite{robinson1987};
and the cost of transport.

\subsection{Torque requirement and actuator sizing}
The assistive torque is the TA moment that the weakness removed, $\tau_\text{req} = 1.2\,\max(0, M_\text{TA,healthy} - M_\text{TA,weak})$, from the stride-averaged moments of the healthy and adapted 25\% gaits; the 20\% is a design margin, and only dorsiflexing moments are kept because that is all TA can do. The device is a brushless DC motor (maxon EC~45 flat, \SI{70}{W}, \SI{24}{V} \cite{maxon411812}), a gear of ratio $N$ and efficiency 0.85, and a series spring $k$ \cite{pratt1995}. For a joint torque $\tau$ and angle $\theta_j$ over one stride, the gear output is $\theta_o = \theta_j + \tau/k$, the motor turns at $N\dot\theta_o$, its torque is $J_m\ddot\theta_m + b_m\dot\theta_m$ plus the reflected load, and $V = Ri + K_e\omega_m$. The design energy per stride is $\int \max(Vi, 0)\,dt$, so the drive does not regenerate (copper losses are inside $Vi$). A design is feasible if peak speed, RMS current (rated \SI{3.21}{A}), peak current (\SI{15}{A}, the drive) and peak voltage stay within limits. The sweep covers $N = 20$ to 150 and 25 stiffnesses from 50 to \SI{2000}{N.m/rad} plus the rigid case. The chosen design is the most compliant one within 2\% of the lowest feasible energy.

The joint trajectory in this calculation is the normative ankle angle \cite{camargo2021} resampled to the model's stride time, not the model's own. The adapted model's ankle slaps down after heel strike (Section~\ref{sec:res-healthy}), reaching \SI{15}{rad/s} and \SI{2100}{rad/s^2}; the reflected rotor inertia then needs more RMS current than the motor's rating for every $N$ and $k$, while the human trajectory peaks at \SI{4}{rad/s} and \SI{66}{rad/s^2}. Both sweeps are reported. Gait curves on the 1\% grid are resampled with periodic cubic splines, because linear interpolation turns the second derivative into impulses at the grid points.

\subsection{Simulated shank gyroscope and event detector}
A Lua script controller reads the sagittal angular velocity of \texttt{tibia\_r} every control step and turns it into a gyroscope signal: samples at \SI{100}{Hz}, white noise of \SI{0.5}{\degree/s} RMS, a bias of \SI{2}{\degree/s}, a \SI{15}{ms} transport delay and a zero-order hold. The causal detector uses the classic shank pattern \cite{aminian2002}: it arms on the positive mid-swing peak, takes heel strike (HS) as the first negative minimum after the terminal-swing zero crossing (feature \emph{min}) or the zero crossing itself (\emph{zero}), and toe-off (TO) as the next negative minimum after a lockout of part of the mean stride. Gait phase is the time since HS divided by the mean of the last three stride times. Each Lua module has a Python twin that reproduces it to rounding error on the same samples, and the noise generator is shared bit for bit. Ground truth is the right vertical ground reaction force crossing 5\% of body weight. Errors are estimate minus truth, latency is confirmation time minus truth, and an estimate pairs with a true event if it falls within $-150$ to $+250$~ms. The detector was validated on three 60~s healthy runs and on the adapted 25\% gait, and then, unchanged, on the shank gyroscopes of the Camargo et al.\ subjects at 1.2~m/s.

\subsection{SEA torque control}
For the chosen $N$ and $k$, the plant is the reflected rotor inertia $J_r = J_m N^2$ and damping $b_r = b_m N^2$ driven by $\eta N K_t i$, the spring torque $\tau = k(\theta_o - \theta_j)$, and the motor winding with a \SI{1}{kHz} PI current loop that saturates at the supply voltage and the drive's peak current. The torque controller runs at \SI{1}{kHz} with one sample of computation delay:
\begin{equation}
u = u_\text{ff} + K_p e + K_i\!\int\! e\,dt - K_d \dot\tau_f,\quad i_\text{cmd} = u / (\eta N K_t),
\end{equation}
with $e = \tau_d - \tau$, the derivative on the measured torque filtered at \SI{100}{Hz}, and anti-windup. The feedforward is the static term $\tau_d$ plus a model term $J_r(\ddot\theta_j + \ddot\tau_d/k) + b_r(\dot\theta_j + \dot\tau_d/k) + K_d\dot\tau_{d,f}$ \cite{vallery2007,paine2014}. The last term matters: without it the derivative on measurement still acts when tracking is perfect, and at the lightly damped locked-output resonance it dominates and opens a 15~dB notch in the closed-loop response. Gains place the locked-output poles at a natural frequency $\omega_n$ with damping 0.7 and the integral zero a factor of five below. The design is the lowest $\omega_n/2\pi$ on a 5~Hz grid, above the locked resonance, whose small-signal response has less than 6~dB of peaking, a bandwidth of at least 10~Hz without the model feedforward, and an RMS error under 10\% of peak on five strides of $\tau_\text{req}$ with saturation. Bode plots come from simulated sine responses with the output locked. The closed loop is then reduced to a pure delay plus a first-order lag, fitted by grid search to the tracking response.

\subsection{Devices and controllers in SCONE}
The device is a moment pair, $+\tau$ on \texttt{talus\_r} and $-\tau$ on \texttt{tibia\_r} (positive dorsiflexes). SCONE external moments persist, so the script adds only the change of moment at each step; a constant-torque pulse test confirmed the bookkeeping. The \emph{passive AFO} is $\tau = -k_\text{AFO}\theta$ with neutral angle 0. The \emph{active AFO} is a finite-state machine driven only by the detector: from HS to 15\% of the estimated cycle, viscous braking $\tau = \max(0, -b\dot\theta)$ with $b = \SI{1}{N.m.s/rad}$; then zero torque until TO; in swing, PD control toward a dorsiflexed target, $\tau = \max(0, k_p(\theta^* - \theta) - k_d\dot\theta)$ with $k_p = \SI{30}{N.m/rad}$ and $k_d = \SI{0.6}{N.m.s/rad}$, until the next HS or a 0.8~s timeout. The gains come from a sizing calculation against the healthy TA moments, not from tuning. The command passes through the fitted actuator model and is clipped at the peak of $\tau_\text{req}$. The device never plantarflexes.

\subsection{Experiment}
Conditions: healthy; drop foot at 25\% TA strength with no device; with the passive AFO; with the active AFO. Each device condition is run with the adapted drop-foot controller \emph{frozen} (the first day with the device) and \emph{re-optimized} for 150 more generations with the device on, from the adapted controller of the same seed (adaptation). Because re-optimization adds generations, the no-device condition is re-optimized for the same 150 generations as a control. The passive stiffness ($k_\text{AFO} = 10$ to \SI{160}{N.m/rad}, factors of two) and the active swing target ($\theta^* = 0$, 5, \SI{10}{\degree}) were each chosen on the frozen controllers by one rule: among the settings with which all seeds walk, the one with toe clearance closest to healthy, the lower setting on a tie. A trial of the pipeline showed that no setting may let all seeds walk, so before the comparison the rule was completed for both devices: among the settings with the most seeds walking, the toe clearance closest to healthy decides, and if no seed walks at all, the longest mean time before the fall.

The hypothesis, committed before any comparison: \emph{at 25\% TA strength, after re-optimization with the device, the active AFO brings mid-swing toe clearance to within 10~mm of healthy without lowering peak ankle push-off power (muscles plus device) by more than 10\% compared with no device, while the best passive AFO lowers it by more than 15\%.} Each of the three parts holds or fails on the seed means; a part that misses by less than one seed SD is inconclusive, and a fall in any seed fails the active parts.

Robustness of the re-optimized conditions was tested at a second speed (all conditions re-optimized with a 1.5~m/s minimum), with the gyroscope noise and delay doubled, with IMU and actuator delays from 0 to 100~ms, and with single \SI{0.2}{s} horizontal pushes on the pelvis as in the SCONE perturbation tutorial, from 25 to 150~N in both directions.
